\section{The literal missing-high interface}
\label{sec:tpc61-physical-interface}

We begin on the same one-side direct slice as TPC--53--60.  The purpose of
this section is to freeze the arithmetic carrier before defining either an
exposure profile or a normalized Gram.  None of the sets below may depend on
the coefficient vector.

\subsection{Scales, rows, and the critical cutoff}

The inherited source row has the unique factorization
\begin{equation}
 m=d_m\ell_m,
 \qquad d_m\asymp D,
 \qquad \ell_m\asymp L,
 \qquad m\asymp Q=DL,
 \label{eq:tpc61-source-scales}
\end{equation}
inside one declared row system.  The orbit time satisfies \(j\asymp J\),
and at the endpoint
\begin{equation}
 \begin{aligned}
 Q&=X^{267/400+o(1)},
 &J&=X^{133/400+o(1)},\\
 D&=X^{10049/52500+o(1)},
 &L&=X^{99979/210000+o(1)}.
 \end{aligned}
 \label{eq:tpc61-endpoint-scales}
\end{equation}
The physical shift \(h_0\ne0\) is fixed.  The equality support is
squarefree and obeys
\begin{equation}
 (d_m,mj+h_0)=(d_m,h_0)=1.
 \label{eq:tpc61-source-coprimality}
\end{equation}

Retain the support endpoints \(D_\pm,M_\pm,J_\pm,N_\pm\) of TPC--59 and
\(\Delta_M=M_+-M_-\).  Throughout, the cutoff \(y\) is
\emph{very-high admissible}:
\begin{equation}
 y>\max\{|h_0|,D_+,J_+,\Delta_M,\sqrt{N_+}\},
 \qquad
 \frac{D_+N_+}{y}<D_-N_-.
 \label{eq:tpc61-very-high-admissible}
\end{equation}
In particular, \(y=C_{\rm vh}Q\) is allowed for a sufficiently large fixed
support-dependent constant \citep{WangTPC59}.  The cutoff is a feasibility
condition, not an energy estimate.

\subsection{Recovered labels and raw amplitudes}

Let \(\cW_H(y)\) be the complete nonzero very-high Walsh carrier on this
slice.  A coordinate \(w\in\cW_H(y)\) uniquely recovers
\begin{equation}
 w=(m,j,\gamma),
 \qquad
 n(w)=mj+h_0,
 \qquad
 p(w)=\Pplus(n(w))>y,
 \label{eq:tpc61-recovered-mjp}
\end{equation}
and, because \(p(w)^2>n(w)\),
\begin{equation}
 r(w)=\frac{n(w)}{p(w)},
 \qquad
 s(w)=d_mr(w).
 \label{eq:tpc61-recovered-rs}
\end{equation}
The native key is
\begin{equation}
 i(w)=(L,\gamma,j)
 \label{eq:tpc61-native-key-left}
\end{equation}
on the displayed left slice.  The right slice exchanges the two row
systems and is added by orthogonal direct sum.  Thus a fixed native key
fixes the side, opposite row, and the same physical orbit time \(j\).

Literal factorization supplies
\begin{equation}
 b_w(\zeta)=\beta_w\zeta_{m(w)},
 \qquad \zeta\in\cK_X,
 \label{eq:tpc61-literal-row-amplitude}
\end{equation}
where \(\beta_w\) retains the source shapes, static mask, opposite-row
factor, orbit multiplier, four-factor profile, and shaped terminal weight,
but not the bounded row coefficient \(\zeta_m\).  It is the stripped
TPC--45 coefficient: the physical high character is applied separately.
No atom-dependent free phase is introduced.

The represented terminal census injects into \(\cW_H(y)\).  Write
\begin{equation}
 \cW_H=\cW_R\,\dot\cup\,\cW_M
 \label{eq:tpc61-RM-partition}
\end{equation}
for its image and complement.  The raw diagonals are
\begin{equation}
 \cD_E(\zeta)
 :=\sum_{w\in\cW_E}|\beta_w|^2|\zeta_{m(w)}|^2,
 \qquad E\in\{R,M,H\},
 \label{eq:tpc61-raw-RMH-diagonals}
\end{equation}
so \(\cD_H=\cD_R+\cD_M\) exactly.

On squarefree support, \(\mu(s)=\mu(d_m)\mu(r)\).  The missing native map
is therefore
\begin{equation}
 (M_X\zeta)_i
 :=-\sum_{\substack{w\in\cW_M\\i(w)=i}}
       \mu(s(w))\beta_w\zeta_{m(w)}.
 \label{eq:tpc61-missing-native-map}
\end{equation}
This is the same sign convention as TPC--58 and TPC--60; inserting another
M\"obius sign would change the physical map \citep{WangTPC58,WangTPC60}.

\subsection{The two singleton incidences}

The very-high inequalities have two consequences used repeatedly.  First,
for fixed \((i,p)\), there is at most one high atom.  Indeed, \(i\) fixes
\(j\); if \(p\) meets rows \(m,m'\), then
\(p\mid(m-m')j\).  The alternative \(p\mid j\) is excluded by
\(p>J_+\), and \(|m-m'|\le\Delta_M<p\) then gives \(m=m'\); unique row
factorization recovers the atom.  Second, for fixed \((i,m)\), the target
\(mj+h_0\), hence \(p,r,s\), is already determined.  Consequently
\begin{equation}
 \#\{w\in\cW_H:(i(w),p(w))=(i,p)\}\le1,
 \qquad
 \#\{w\in\cW_H:(i(w),m(w))=(i,m)\}\le1.
 \label{eq:tpc61-ip-im-singletons}
\end{equation}
These are coordinate incidences, not statements that a native output has
only one prime or only one row.

\subsection{When missing means cofactor mismatch}

The exposure theorem below uses the following explicit specialization.
\begin{definition}[Identical-support specialization]
\label{def:tpc61-identical-support}
The Walsh and terminal constructions have identical row, orbit,
opposite-row, terminal-prime, literal-mask, and profile supports, with no
additional terminal-only pruning.  If
\begin{equation}
 \cI_m:=\{r:(m,r)\in\cI_X\},
 \label{eq:tpc61-retained-cofactor-row}
\end{equation}
then
\begin{equation}
 w\in\cW_M
 \quad\Longleftrightarrow\quad
 r(w)\notin\cI_{m(w)}.
 \label{eq:tpc61-missing-iff-cofactor}
\end{equation}
\end{definition}
TPC--60 proves this conditional support identity.  Outside this
specialization one must retain the full first-failure partition; the
cofactor complement alone need not equal \(\cW_M\).

\subsection{Six-item physical audit}

The construction satisfies the required interface audit.
\begin{enumerate}[leftmargin=2.3em]
 \item Every coefficient is the literal \(\beta_w\zeta_{m(w)}\).
 \item Every target is \(mj+h_0\) for one fixed nonzero physical shift.
 \item All atomic energies use raw counting measure, with no fiber average.
 \item The row, largest prime, cofactor, residual label, and native key are
       canonically recovered; only identically inactive rows are removed.
 \item The carriers \(\cW_H,\cW_R,\cW_M\) and \(\cI_X\) are fixed before
       \(\zeta\) is chosen.
 \item Exposure and Gram identities have zero structural cost, but every
       eventual reassembly loss remains in the strict \(1/400\) ledger.
\end{enumerate}
